\section{Проектирование архитектуры прототипа системы профилирования}
\label{sec:razdel3}

\subsection{Обоснование выбора модели жизненного цикла и инструментальных средств}
\label{sec:3.1}

Работа носит исследовательско-прикладной характер: на момент постановки
задачи (раздел~\ref{sec:razdel1}) не было известно заранее, какой именно
набор параметров критерия~\eqref{eq:criterion} ($T$, $\alpha$,
\texttt{-count}) окажется работоспособным на выбранной тестовой нагрузке —
ответ на этот вопрос сам является результатом эксперимента
(раздел~\ref{sec:razdel5}). По этой причине для разработки прототипа
выбрана \emph{эволюционная (итеративная) модель жизненного цикла} вместо
каскадной: каждая итерация добавляла минимально достаточный слой
функциональности (сначала тестовое приложение и бенчмарки, затем
статистическое сравнение, затем pprof-диагностика, затем эксперименты) с
немедленной проверкой работоспособности перед переходом к следующему
слою. Такой подход соответствует рекомендациям ГОСТ Р ИСО/МЭК
12207~\cite{gost12207} в части организации процессов жизненного цикла
программных средств для случая, когда требования уточняются по ходу
разработки; в отличие от стадийной модели ГОСТ 34.601-90~\cite{gost34601}
(ориентированной на последовательные стадии "техническое задание — эскизный
проект — технический проект — рабочая документация"), выбранная модель не
предполагает фиксации всех проектных решений до начала реализации. Такой
подход согласуется с принципом коротких обратных связей, описанным в
контексте практик непрерывной интеграции у Fowler~\cite{fowler-ci} и
Humble \& Farley~\cite{humble-farley2010}.

Выбор инструментальных средств обоснован следующими требованиями:

\begin{itemize}
\item \textbf{Язык реализации — Go 1.22~\cite{go-release-notes}.} Определён
темой работы; кроме того, Go — единственный из мейнстримных языков, поставляющий
статистический сравнитель бенчмарков (\texttt{benchstat}) и профилировщик
(\texttt{pprof}) как часть официальной экосистемы, а не сторонний пакет,
что снижает поверхность интеграции.
\item \textbf{CI-платформа — GitHub Actions.} Выбрана как наиболее
распространённая платформа для open-source Go-проектов на момент
разработки, с нативной поддержкой шага \texttt{\$GITHUB\_STEP\_SUMMARY} для
публикации markdown-отчёта прямо в интерфейсе pull request и встроенным
механизмом артефактов (\texttt{actions/\allowbreak upload-artifact}) для сохранения
бинарных pprof-профилей~\cite{gh-actions-doc}.
\item \textbf{Формат хранения эталона.} Агрегированный эталон —
текстовый вывод \texttt{go test -bench} (человекочитаемый, диффится в git);
профильный эталон — бинарные \texttt{.pprof}-файлы (protobuf/gzip),
хранящиеся в репозитории как обычные версионируемые бинарные артефакты, без
внешнего объектного хранилища — оправдано малым размером тестового
приложения (профили десятки — сотни килобайт).
\item \textbf{Язык скриптов интеграции — POSIX-совместимый \texttt{bash}.}
Выбран вместо, например, Go-программы-оркестратора, поскольку вся логика
сводится к последовательному вызову внешних инструментов (\texttt{go
test}, \texttt{benchstat}, \texttt{go tool pprof}) и разбору их текстового
вывода — задача, для которой оболочка является более простым и прозрачным
решением, чем компилируемая программа с дополнительной точкой сборки.
\end{itemize}

Требования к прототипу как программной системе сформулированы следующим
образом. \textbf{Функциональные:} (1) выполнить набор бенчмарков заданное
число раз; (2) статистически сравнить результат с эталоном по
критерию~\eqref{eq:criterion} и вернуть код завершения, пригодный для
блокировки слияния; (3) снять CPU- и heap-профили того же прогона и
сравнить их с эталонными профилями на уровне отдельных функций;
(4) опубликовать оба результата в человекочитаемом виде в теле pull
request. \textbf{Системные:} воспроизводимость (одинаковые входные данные
и параметры дают одинаковое решение с точностью до статистического шума
запуска), независимость от внешнего состояния (отсутствие сетевых
зависимостей во время сравнения, кроме однократной установки
\texttt{benchstat}), совместимость с \texttt{ubuntu-latest}-раннером
GitHub Actions без дополнительных прав. \textbf{Требования к
пользовательскому интерфейсу:} система не имеет графического интерфейса —
её "интерфейсом" является (а) консольный вывод скриптов при локальном
запуске и (б) markdown-отчёт в \texttt{\$GITHUB\_STEP\_SUMMARY} при запуске
в CI; оба должны оставаться читаемыми без дополнительной постобработки
(табличный markdown, без ANSI-последовательностей).

\subsection{Архитектура прототипа}
\label{sec:3.2}

Архитектура прототипа (рис.~\ref{fig:3.1}) состоит из четырёх слоёв:
приложения под нагрузкой (\texttt{app/\allowbreak }), эталонных данных
(\texttt{benchmarks/\allowbreak }), слоя сравнения (\texttt{scripts/\allowbreak }) и слоя оркестрации
CI (\texttt{.github/\allowbreak workflows/\allowbreak }). Слой сравнения намеренно отделён от
приложения: ни один из скриптов сравнения не импортирует пакет
\texttt{app} и не содержит специфичной для него логики — единственная связь
между слоями осуществляется через текстовый/бинарный формат вывода
\texttt{go test}, что позволяет использовать те же скрипты для любого
Go-пакета с бенчмарками без модификации (общность подтверждена тем, что
\texttt{scripts/\allowbreak compare.sh} параметризован переменной окружения
\texttt{PKG}).

\begin{figure}[htbp]
\centering
\resizebox{\linewidth}{!}{%
\begin{tikzpicture}[
  node distance=7mm and 10mm,
  layer/.style={draw, rounded corners, minimum height=9mm, minimum width=34mm, align=center, font=\small},
  group/.style={draw, dashed, rounded corners, inner sep=6mm, font=\bfseries\small},
  arr/.style={-{Latex[length=2mm]}}
]
\node[layer] (proc) {app/processor.go\\(6 функций бизнес-логики)};
\node[layer, below=of proc] (bench) {app/bench\_test.go\\(6 бенчмарков)};
\node[layer, below=of bench] (unit) {app/processor\_test.go\\(unit-тесты)};

\node[layer, right=30mm of proc] (basetxt) {benchmarks/baseline.txt};
\node[layer, below=of basetxt] (basepprof) {benchmarks/baseline-*.pprof};

\node[layer, right=30mm of basetxt] (lib) {scripts/lib.sh\\(разбор CSV, критерий)};
\node[layer, below=of lib] (compare) {scripts/compare.sh};
\node[layer, below=of compare] (diff) {scripts/profile-diff.sh};
\node[layer, below=of diff] (exp) {scripts/experiments/*.sh};

\node[layer, right=30mm of lib] (ci) {.github/workflows/\\perf-check.yml};

\begin{scope}[on background layer]
\node[group, fit=(proc)(bench)(unit), label={[font=\small]above:Приложение под нагрузкой}] {};
\node[group, fit=(basetxt)(basepprof), label={[font=\small]above:Эталонные данные}] {};
\node[group, fit=(lib)(compare)(diff)(exp), label={[font=\small]above:Слой сравнения}] {};
\node[group, fit=(ci), label={[font=\small]above:Оркестрация CI}] {};
\end{scope}

\draw[arr] (bench) -- (compare);
\draw[arr] (bench) -- (diff);
\draw[arr] (basetxt) -- (compare);
\draw[arr] (basepprof) -- (diff);
\draw[arr] (lib) -- (compare);
\draw[arr] (lib) -- (exp);
\draw[arr] (ci) -- (compare);
\draw[arr] (ci) -- (diff);
\end{tikzpicture}%
}
\caption{Архитектура прототипа системы непрерывного профилирования}
\label{fig:3.1}
\end{figure}

\subsection{Состав и структура основных компонентов}
\label{sec:3.3}

\textbf{Приложение под нагрузкой (\texttt{app/\allowbreak }).} Не является предметом
исследования само по себе — это управляемая, воспроизводимая нагрузка,
на которой демонстрируется работа системы профилирования. Состоит из
шести функций обработки записей (\texttt{ParseRecords}, \texttt{FilterActive},
\texttt{FindByID}, \texttt{Aggregate}, \texttt{FormatNames},
\texttt{Deduplicate}), каждая из которых сознательно выбрана так, чтобы
охватить разный характер нагрузки: разбор JSON (I/O-подобная, много
аллокаций), линейная фильтрация и поиск (CPU-bound без аллокаций),
агрегация в map (смешанная), дедупликация (структура данных,
чувствительная к алгоритмической сложности — используется как раз для
сценария \texttt{regression/\allowbreak quadratic-dedup} в разделе~\ref{sec:razdel5}).

\textbf{Эталонные данные (\texttt{benchmarks/\allowbreak }).} Два независимых
представления одного и того же эталонного прогона: \texttt{baseline.txt} —
агрегат для \texttt{benchstat}, и пара \texttt{baseline-cpu.pprof} /
\texttt{baseline-mem.pprof} — профили для \texttt{go tool pprof
-diff\_base}. Оба генерируются отдельными идемпотентными скриптами
(\texttt{scripts/\allowbreak compare.sh} неявно ожидает первый, явный генератор для
профилей — \texttt{scripts/\allowbreak capture-baseline-profiles.sh}), обновляются
разработчиком осознанно при намеренном изменении производительности
эталона, а не автоматически при каждом коммите.

\textbf{Слой сравнения (\texttt{scripts/\allowbreak }).} Реализует оба уровня
критерия из раздела~\ref{sec:razdel2}: \texttt{lib.sh} содержит общую
функцию разбора CSV-отчёта \texttt{benchstat} и вычисления
критерия~\eqref{eq:criterion} (переиспользуется тремя вызывающими
скриптами, что исключает рассинхронизацию логики критерия);
\texttt{compare.sh} — решающий шаг (раздел~\ref{sec:2.2}); реализация
подробно описана в разделе~\ref{sec:razdel4}); \texttt{profile-diff.sh} —
диагностический шаг (раздел~\ref{sec:2.1}); \texttt{scripts/\allowbreak experiments/\allowbreak }
— три скрипта, реализующие эксперименты раздела~\ref{sec:razdel5}
(чувствительность порога, накладные расходы) как переиспользование тех же
примитивов сравнения в цикле по сетке параметров.

\textbf{Оркестрация CI (\texttt{.github/\allowbreak workflows/\allowbreak perf-check.yml}).}
Единственный workflow, триггер — \texttt{pull\_request} в \texttt{main}.
Запускает решающий и диагностический шаги последовательно в одной джобе на
\texttt{ubuntu-latest}, публикует markdown-отчёт в
\texttt{\$GITHUB\_STEP\_SUMMARY} и сохраняет pprof-профили текущего
прогона как артефакт джобы (\texttt{actions/\allowbreak upload-artifact}) — что даёт
разработчику возможность локально скачать профиль и открыть его
интерактивным веб-интерфейсом (\texttt{go tool pprof -http}) без
повторного запуска бенчмарков.

\subsection{Выводы}

\begin{enumerate}
\item Обоснован выбор эволюционной (итеративной) модели жизненного цикла
для разработки прототипа как более соответствующей исследовательскому
характеру задачи по сравнению с каскадной моделью, со ссылкой на
ГОСТ Р ИСО/МЭК 12207~\cite{gost12207}.
\item Обоснован выбор инструментальных средств: Go 1.22 (единственная
экосистема с официальным статистическим сравнителем бенчмарков и
профилировщиком), GitHub Actions (нативная поддержка markdown-отчётов в PR
и артефактов), POSIX \texttt{bash} для слоя интеграции.
\item Сформулированы функциональные, системные требования и требования к
интерфейсу прототипа — показано, что при отсутствии графического
интерфейса роль пользовательского интерфейса выполняет
структурированный (markdown, табличный) консольный/CI-вывод.
\item Разработана архитектура прототипа (рис.~\ref{fig:3.1}) из четырёх
слабо связанных слоёв (приложение, эталонные данные, слой сравнения,
оркестрация CI), связь между которыми ограничена текстовым/бинарным
форматом вывода \texttt{go test}, что делает слой сравнения переиспользуемым
для произвольного Go-пакета с бенчмарками.
\item Описан состав и назначение шести основных компонентов системы
(\texttt{app/\allowbreak }, \texttt{benchmarks/\allowbreak }, \texttt{scripts/\allowbreak lib.sh},
\texttt{scripts/\allowbreak compare.sh}, \texttt{scripts/\allowbreak profile-diff.sh},
\texttt{.github/\allowbreak workflows/\allowbreak perf-check.yml}); показано, что оба уровня
критерия (раздел~\ref{sec:razdel2}) реализованы как независимые шаги одной
CI-джобы, что соответствует модели процесса, разработанной в
разделе~\ref{sec:2.1}.
\end{enumerate}
